%!TEX root = ../template.tex
%%%%%%%%%%%%%%%%%%%%%%%%%%%%%%%%%%%%%%%%%%%%%%%%%%%%%%%%%%%%%%%%%%%%
%% symbols.tex
%% NOVA thesis document file
%%
%% Symbols definition
%%%%%%%%%%%%%%%%%%%%%%%%%%%%%%%%%%%%%%%%%%%%%%%%%%%%%%%%%%%%%%%%%%%%
\typeout{NT FILE symbols.tex}%

% [REV-2026-09-13][FR-02][P3][NOTACAO] Definir a notação física e estatística com unidades
% Trocar símbolos de exemplo por E_t, Q_t, Z_t, a, b, e_t, sigma residual,
% sigma da inovação, rho, delta_MAE, cobertura, k e h do CUSUM.
% Dar unidades e separar índice temporal de ano. A unidade do declive
% com E em t GNE/d e Q em kt/d é t GNE/kt: o dia cancela.
% Distinguir IC de parâmetro/média, intervalo preditivo e banda de alarme
% no Apêndice G e nas legendas.
% [/REV-2026-09-13]
\glsxtrnewsymbol[sort=mu,description={Mu}]{mu}{\ensuremath{\mu}}
\glsxtrnewsymbol[description={the numerical value of pi}]{pi}{\ensuremath{\pi}}
\glsxtrnewsymbol[description={the radius of a circle}]{r}{\ensuremath{r}}
% \glsxtrnewsymbol[description={Total Angular Momentum}]{jay}{\ensuremath{\vv{J}}}
